% LingBot-VLA 2.0 4090 集群交接文档(2026-08-15 更新)
% 编译:tectonic -X compile lingbot_vla_v2_4090_handoff.tex
\documentclass[11pt]{ctexart}
% cnreport-preamble.tex — 中文技术报告导言区(Cosmos/arXiv 单栏风格)
% 用法: \documentclass[11pt]{ctexart} 之后 % cnreport-preamble.tex — 中文技术报告导言区(Cosmos/arXiv 单栏风格)
% 用法: \documentclass[11pt]{ctexart} 之后 % cnreport-preamble.tex — 中文技术报告导言区(Cosmos/arXiv 单栏风格)
% 用法: \documentclass[11pt]{ctexart} 之后 \input{cnreport-preamble}
% 编译: tectonic(XeTeX 引擎),需要系统安装 Noto Serif/Sans CJK SC 字体
\usepackage[a4paper,top=2.7cm,bottom=2.7cm,left=2.5cm,right=2.5cm,headheight=15pt]{geometry}
\usepackage{fontspec}
% 正文:拉丁衬线 + 思源宋体;标题:拉丁无衬线 + 思源黑体
% 按文件名从 TeX 字体树加载(tectonic bundle 自带 TeX Gyre;系统字体需 Noto CJK SC)
\setmainfont{texgyretermes}[Extension=.otf,UprightFont=*-regular,BoldFont=*-bold,ItalicFont=*-italic,BoldItalicFont=*-bolditalic]
\setsansfont{texgyreheros}[Extension=.otf,UprightFont=*-regular,BoldFont=*-bold,ItalicFont=*-italic,BoldItalicFont=*-bolditalic]
\setmonofont{texgyrecursor}[Extension=.otf,UprightFont=*-regular,BoldFont=*-bold,Scale=0.88]
\setCJKmainfont{Noto Serif CJK SC}[BoldFont={Noto Serif CJK SC Bold}]
\setCJKsansfont{Noto Sans CJK SC}[BoldFont={Noto Sans CJK SC Bold}]
\setCJKmonofont{Noto Sans Mono CJK SC}
% 带圈数字 ①-⑳ 等划入 CJK 类,用 CJK 字体渲染(xeCJK 默认按拉丁处理)
\xeCJKDeclareCharClass{CJK}{"2460 -> "24FF}

\usepackage{amsmath,amssymb,bm}
\usepackage{booktabs}
\usepackage{graphicx}
\usepackage{xcolor}
\usepackage{titlesec}
\usepackage{fancyhdr}
\usepackage{enumitem}
\usepackage{caption}
\usepackage{listings}

% ---- 配色(主色可改) ----
\definecolor{accent}{HTML}{0E5FA8}     % 强调蓝(链接/节号)
\definecolor{good}{HTML}{0E7A4D}       % 正向结果
\definecolor{warn}{HTML}{B45309}       % 负向/警示
\definecolor{faint}{HTML}{6B7889}      % 弱化文字
\definecolor{codebg}{HTML}{F0F2F5}
\definecolor{rule}{HTML}{C9CFD8}

\usepackage[colorlinks=true,linkcolor=accent,citecolor=accent,urlcolor=accent]{hyperref}

% ---- 章节标题:无衬线粗体,节号着色 ----
\titleformat{\section}{\Large\bfseries\sffamily}{\textcolor{accent}{\thesection}}{0.6em}{}
\titleformat{\subsection}{\large\bfseries\sffamily}{\textcolor{accent}{\thesubsection}}{0.5em}{}
\titleformat{\subsubsection}{\normalsize\bfseries\sffamily}{\textcolor{accent}{\thesubsubsection}}{0.5em}{}
\titlespacing{\section}{0pt}{1.6em}{0.7em}
\titlespacing{\subsection}{0pt}{1.2em}{0.5em}

% ---- 页眉页脚(Cosmos 式:左标题右日期,下细线) ----
\newcommand{\runningtitle}{}  % 主文档用 \renewcommand 设置
\newcommand{\reportdate}{}
\newcommand{\reportfooter}{}
\fancypagestyle{cnreport}{%
  \fancyhf{}
  \fancyhead[L]{\small\sffamily \runningtitle}
  \fancyhead[R]{\small\sffamily \reportdate}
  \fancyfoot[L]{\small\color{faint}\sffamily \reportfooter}
  \fancyfoot[R]{\small\sffamily \thepage}
  \renewcommand{\headrulewidth}{0.4pt}
  \renewcommand{\footrulewidth}{0pt}
}
\pagestyle{cnreport}

% ---- 题注与列表 ----
\captionsetup{font=small,labelfont={bf,sf},skip=6pt}
\setlist{itemsep=2pt,topsep=4pt,parsep=0pt}

% ---- 代码 ----
\lstset{
  basicstyle=\ttfamily\small,
  backgroundcolor=\color{codebg},
  frame=single, framerule=0.3pt, rulecolor=\color{rule},
  xleftmargin=6pt, xrightmargin=6pt,
  breaklines=true, columns=fullflexible, keepspaces=true,
  aboveskip=8pt, belowskip=8pt,
}

% ---- 报告头(kicker / 标题 / 副题 / 署名 / 摘要 / 关键词) ----
% 用法: \reporthead{kicker-left}{kicker-right}{标题}{副题}{署名行}
\newcommand{\reporthead}[5]{%
  \begin{center}
  {\small\sffamily\bfseries\color{accent} #1 \hfill #2}\\[2pt]
  {\color{black}\rule{\linewidth}{1.2pt}}\\[14pt]
  {\LARGE\sffamily\bfseries #3\par}\vspace{8pt}
  {\large\color{faint} #4\par}\vspace{10pt}
  {\small\sffamily\color{faint} #5\par}\vspace{4pt}
  {\color{rule}\rule{\linewidth}{0.4pt}}
  \end{center}\vspace{6pt}
}
% 摘要块: \begin{cnabstract} ... \end{cnabstract}
\newenvironment{cnabstract}{%
  \begin{quote}\small\color{faint}\sffamily\bfseries\color{black} 摘要 \quad\normalfont\rmfamily\color{faint!80!black}%
}{\end{quote}\vspace{2pt}}
% 重点结论框: \begin{quotebox} ... \end{quotebox}
\newenvironment{quotebox}{%
  \begin{quote}\small
  {\color{rule}\rule{\linewidth}{0.8pt}}\par\vspace{2pt}%
}{\par\vspace{2pt}{\color{rule}\rule{\linewidth}{0.8pt}}\end{quote}\vspace{2pt}}

\newcommand{\keywords}[1]{\begin{quote}\small\color{faint}{\bfseries\sffamily\color{black} 关键词:\ } #1\end{quote}}

% 表内彩色标记
\newcommand{\upgood}[1]{\textcolor{good}{\bfseries #1}}
\newcommand{\downbad}[1]{\textcolor{warn}{\bfseries #1}}
\newcommand{\dimtext}[1]{\textcolor{faint}{#1}}

\setlength{\parskip}{2pt}
\linespread{1.12}
\emergencystretch=3em  % CJK 段落容忍度,消除 overfull

% 编译: tectonic(XeTeX 引擎),需要系统安装 Noto Serif/Sans CJK SC 字体
\usepackage[a4paper,top=2.7cm,bottom=2.7cm,left=2.5cm,right=2.5cm,headheight=15pt]{geometry}
\usepackage{fontspec}
% 正文:拉丁衬线 + 思源宋体;标题:拉丁无衬线 + 思源黑体
% 按文件名从 TeX 字体树加载(tectonic bundle 自带 TeX Gyre;系统字体需 Noto CJK SC)
\setmainfont{texgyretermes}[Extension=.otf,UprightFont=*-regular,BoldFont=*-bold,ItalicFont=*-italic,BoldItalicFont=*-bolditalic]
\setsansfont{texgyreheros}[Extension=.otf,UprightFont=*-regular,BoldFont=*-bold,ItalicFont=*-italic,BoldItalicFont=*-bolditalic]
\setmonofont{texgyrecursor}[Extension=.otf,UprightFont=*-regular,BoldFont=*-bold,Scale=0.88]
\setCJKmainfont{Noto Serif CJK SC}[BoldFont={Noto Serif CJK SC Bold}]
\setCJKsansfont{Noto Sans CJK SC}[BoldFont={Noto Sans CJK SC Bold}]
\setCJKmonofont{Noto Sans Mono CJK SC}
% 带圈数字 ①-⑳ 等划入 CJK 类,用 CJK 字体渲染(xeCJK 默认按拉丁处理)
\xeCJKDeclareCharClass{CJK}{"2460 -> "24FF}

\usepackage{amsmath,amssymb,bm}
\usepackage{booktabs}
\usepackage{graphicx}
\usepackage{xcolor}
\usepackage{titlesec}
\usepackage{fancyhdr}
\usepackage{enumitem}
\usepackage{caption}
\usepackage{listings}

% ---- 配色(主色可改) ----
\definecolor{accent}{HTML}{0E5FA8}     % 强调蓝(链接/节号)
\definecolor{good}{HTML}{0E7A4D}       % 正向结果
\definecolor{warn}{HTML}{B45309}       % 负向/警示
\definecolor{faint}{HTML}{6B7889}      % 弱化文字
\definecolor{codebg}{HTML}{F0F2F5}
\definecolor{rule}{HTML}{C9CFD8}

\usepackage[colorlinks=true,linkcolor=accent,citecolor=accent,urlcolor=accent]{hyperref}

% ---- 章节标题:无衬线粗体,节号着色 ----
\titleformat{\section}{\Large\bfseries\sffamily}{\textcolor{accent}{\thesection}}{0.6em}{}
\titleformat{\subsection}{\large\bfseries\sffamily}{\textcolor{accent}{\thesubsection}}{0.5em}{}
\titleformat{\subsubsection}{\normalsize\bfseries\sffamily}{\textcolor{accent}{\thesubsubsection}}{0.5em}{}
\titlespacing{\section}{0pt}{1.6em}{0.7em}
\titlespacing{\subsection}{0pt}{1.2em}{0.5em}

% ---- 页眉页脚(Cosmos 式:左标题右日期,下细线) ----
\newcommand{\runningtitle}{}  % 主文档用 \renewcommand 设置
\newcommand{\reportdate}{}
\newcommand{\reportfooter}{}
\fancypagestyle{cnreport}{%
  \fancyhf{}
  \fancyhead[L]{\small\sffamily \runningtitle}
  \fancyhead[R]{\small\sffamily \reportdate}
  \fancyfoot[L]{\small\color{faint}\sffamily \reportfooter}
  \fancyfoot[R]{\small\sffamily \thepage}
  \renewcommand{\headrulewidth}{0.4pt}
  \renewcommand{\footrulewidth}{0pt}
}
\pagestyle{cnreport}

% ---- 题注与列表 ----
\captionsetup{font=small,labelfont={bf,sf},skip=6pt}
\setlist{itemsep=2pt,topsep=4pt,parsep=0pt}

% ---- 代码 ----
\lstset{
  basicstyle=\ttfamily\small,
  backgroundcolor=\color{codebg},
  frame=single, framerule=0.3pt, rulecolor=\color{rule},
  xleftmargin=6pt, xrightmargin=6pt,
  breaklines=true, columns=fullflexible, keepspaces=true,
  aboveskip=8pt, belowskip=8pt,
}

% ---- 报告头(kicker / 标题 / 副题 / 署名 / 摘要 / 关键词) ----
% 用法: \reporthead{kicker-left}{kicker-right}{标题}{副题}{署名行}
\newcommand{\reporthead}[5]{%
  \begin{center}
  {\small\sffamily\bfseries\color{accent} #1 \hfill #2}\\[2pt]
  {\color{black}\rule{\linewidth}{1.2pt}}\\[14pt]
  {\LARGE\sffamily\bfseries #3\par}\vspace{8pt}
  {\large\color{faint} #4\par}\vspace{10pt}
  {\small\sffamily\color{faint} #5\par}\vspace{4pt}
  {\color{rule}\rule{\linewidth}{0.4pt}}
  \end{center}\vspace{6pt}
}
% 摘要块: \begin{cnabstract} ... \end{cnabstract}
\newenvironment{cnabstract}{%
  \begin{quote}\small\color{faint}\sffamily\bfseries\color{black} 摘要 \quad\normalfont\rmfamily\color{faint!80!black}%
}{\end{quote}\vspace{2pt}}
% 重点结论框: \begin{quotebox} ... \end{quotebox}
\newenvironment{quotebox}{%
  \begin{quote}\small
  {\color{rule}\rule{\linewidth}{0.8pt}}\par\vspace{2pt}%
}{\par\vspace{2pt}{\color{rule}\rule{\linewidth}{0.8pt}}\end{quote}\vspace{2pt}}

\newcommand{\keywords}[1]{\begin{quote}\small\color{faint}{\bfseries\sffamily\color{black} 关键词:\ } #1\end{quote}}

% 表内彩色标记
\newcommand{\upgood}[1]{\textcolor{good}{\bfseries #1}}
\newcommand{\downbad}[1]{\textcolor{warn}{\bfseries #1}}
\newcommand{\dimtext}[1]{\textcolor{faint}{#1}}

\setlength{\parskip}{2pt}
\linespread{1.12}
\emergencystretch=3em  % CJK 段落容忍度,消除 overfull

% 编译: tectonic(XeTeX 引擎),需要系统安装 Noto Serif/Sans CJK SC 字体
\usepackage[a4paper,top=2.7cm,bottom=2.7cm,left=2.5cm,right=2.5cm,headheight=15pt]{geometry}
\usepackage{fontspec}
% 正文:拉丁衬线 + 思源宋体;标题:拉丁无衬线 + 思源黑体
% 按文件名从 TeX 字体树加载(tectonic bundle 自带 TeX Gyre;系统字体需 Noto CJK SC)
\setmainfont{texgyretermes}[Extension=.otf,UprightFont=*-regular,BoldFont=*-bold,ItalicFont=*-italic,BoldItalicFont=*-bolditalic]
\setsansfont{texgyreheros}[Extension=.otf,UprightFont=*-regular,BoldFont=*-bold,ItalicFont=*-italic,BoldItalicFont=*-bolditalic]
\setmonofont{texgyrecursor}[Extension=.otf,UprightFont=*-regular,BoldFont=*-bold,Scale=0.88]
\setCJKmainfont{Noto Serif CJK SC}[BoldFont={Noto Serif CJK SC Bold}]
\setCJKsansfont{Noto Sans CJK SC}[BoldFont={Noto Sans CJK SC Bold}]
\setCJKmonofont{Noto Sans Mono CJK SC}
% 带圈数字 ①-⑳ 等划入 CJK 类,用 CJK 字体渲染(xeCJK 默认按拉丁处理)
\xeCJKDeclareCharClass{CJK}{"2460 -> "24FF}

\usepackage{amsmath,amssymb,bm}
\usepackage{booktabs}
\usepackage{graphicx}
\usepackage{xcolor}
\usepackage{titlesec}
\usepackage{fancyhdr}
\usepackage{enumitem}
\usepackage{caption}
\usepackage{listings}

% ---- 配色(主色可改) ----
\definecolor{accent}{HTML}{0E5FA8}     % 强调蓝(链接/节号)
\definecolor{good}{HTML}{0E7A4D}       % 正向结果
\definecolor{warn}{HTML}{B45309}       % 负向/警示
\definecolor{faint}{HTML}{6B7889}      % 弱化文字
\definecolor{codebg}{HTML}{F0F2F5}
\definecolor{rule}{HTML}{C9CFD8}

\usepackage[colorlinks=true,linkcolor=accent,citecolor=accent,urlcolor=accent]{hyperref}

% ---- 章节标题:无衬线粗体,节号着色 ----
\titleformat{\section}{\Large\bfseries\sffamily}{\textcolor{accent}{\thesection}}{0.6em}{}
\titleformat{\subsection}{\large\bfseries\sffamily}{\textcolor{accent}{\thesubsection}}{0.5em}{}
\titleformat{\subsubsection}{\normalsize\bfseries\sffamily}{\textcolor{accent}{\thesubsubsection}}{0.5em}{}
\titlespacing{\section}{0pt}{1.6em}{0.7em}
\titlespacing{\subsection}{0pt}{1.2em}{0.5em}

% ---- 页眉页脚(Cosmos 式:左标题右日期,下细线) ----
\newcommand{\runningtitle}{}  % 主文档用 \renewcommand 设置
\newcommand{\reportdate}{}
\newcommand{\reportfooter}{}
\fancypagestyle{cnreport}{%
  \fancyhf{}
  \fancyhead[L]{\small\sffamily \runningtitle}
  \fancyhead[R]{\small\sffamily \reportdate}
  \fancyfoot[L]{\small\color{faint}\sffamily \reportfooter}
  \fancyfoot[R]{\small\sffamily \thepage}
  \renewcommand{\headrulewidth}{0.4pt}
  \renewcommand{\footrulewidth}{0pt}
}
\pagestyle{cnreport}

% ---- 题注与列表 ----
\captionsetup{font=small,labelfont={bf,sf},skip=6pt}
\setlist{itemsep=2pt,topsep=4pt,parsep=0pt}

% ---- 代码 ----
\lstset{
  basicstyle=\ttfamily\small,
  backgroundcolor=\color{codebg},
  frame=single, framerule=0.3pt, rulecolor=\color{rule},
  xleftmargin=6pt, xrightmargin=6pt,
  breaklines=true, columns=fullflexible, keepspaces=true,
  aboveskip=8pt, belowskip=8pt,
}

% ---- 报告头(kicker / 标题 / 副题 / 署名 / 摘要 / 关键词) ----
% 用法: \reporthead{kicker-left}{kicker-right}{标题}{副题}{署名行}
\newcommand{\reporthead}[5]{%
  \begin{center}
  {\small\sffamily\bfseries\color{accent} #1 \hfill #2}\\[2pt]
  {\color{black}\rule{\linewidth}{1.2pt}}\\[14pt]
  {\LARGE\sffamily\bfseries #3\par}\vspace{8pt}
  {\large\color{faint} #4\par}\vspace{10pt}
  {\small\sffamily\color{faint} #5\par}\vspace{4pt}
  {\color{rule}\rule{\linewidth}{0.4pt}}
  \end{center}\vspace{6pt}
}
% 摘要块: \begin{cnabstract} ... \end{cnabstract}
\newenvironment{cnabstract}{%
  \begin{quote}\small\color{faint}\sffamily\bfseries\color{black} 摘要 \quad\normalfont\rmfamily\color{faint!80!black}%
}{\end{quote}\vspace{2pt}}
% 重点结论框: \begin{quotebox} ... \end{quotebox}
\newenvironment{quotebox}{%
  \begin{quote}\small
  {\color{rule}\rule{\linewidth}{0.8pt}}\par\vspace{2pt}%
}{\par\vspace{2pt}{\color{rule}\rule{\linewidth}{0.8pt}}\end{quote}\vspace{2pt}}

\newcommand{\keywords}[1]{\begin{quote}\small\color{faint}{\bfseries\sffamily\color{black} 关键词:\ } #1\end{quote}}

% 表内彩色标记
\newcommand{\upgood}[1]{\textcolor{good}{\bfseries #1}}
\newcommand{\downbad}[1]{\textcolor{warn}{\bfseries #1}}
\newcommand{\dimtext}[1]{\textcolor{faint}{#1}}

\setlength{\parskip}{2pt}
\linespread{1.12}
\emergencystretch=3em  % CJK 段落容忍度,消除 overfull


\renewcommand{\runningtitle}{LingBot-VLA 2.0 4090 集群交接}
\renewcommand{\reportdate}{2026 年 8 月 15 日}
\renewcommand{\reportfooter}{accel-bench · TR-2026-08-15-HANDOFF}

\begin{document}

\reporthead{交接文档 · TR-2026-08-15-HANDOFF}{2026 年 8 月 15 日}
{LingBot-VLA 2.0 推理加速:4090 集群实验全录与收官状态}
{机器现状 / 代码地图 / 实验台账 / CUDA graph 收官与训练打通}
{平台:4$\times$RTX 4090 24GB(compshare)\quad 代码:lerobot pr3967-combined + fork miracle-techlink}

\begin{quotebox}
\textbf{一句话现状:}路线 A(CUDA graph 去噪环整环捕获)\textbf{已打通}并四路逐位验证全零
(commit \texttt{df0cf53e},库级开关 \texttt{use\_cudagraph\_denoise})。最终 bake 配置
(eager attn + gridthw + cudagraph + \texttt{num\_steps=7})库级严格口径 \textbf{130.9ms}、
4 步 \textbf{108.5ms}——与官方 README 的 130ms claim 数字持平,但\textbf{公开可复现}
(官方 claim 经两路独立调查判定不可复现,见 §3.3)。调度层 RTC 三臂实测完成:
有效反应延迟 sync 1.0--1.2s $\to$ async $\approx$700ms $\to$ RTC $\approx$580ms。
训练侧:单卡 24GB(gradient\_checkpointing,峰值 22.0GB)与 2$\times$24GB FSDP2
(四件套,$\approx$23s/步)均已打通。
\end{quotebox}

\section{4090 机器现状}

\subsection{硬件与环境(含重要更正)}
\begin{itemize}
\item \textbf{更更正:该机是 4$\times$ 满血 RTX 4090 (sm\_89, Ada),不是 4090D,也不是
sm\_120}。\texttt{torch.cuda.}\texttt{get\_device\_capability} 实测 (8,9)。此前文档
(含 pipeline PDF)中 ``4090D/sm\_120'' 的表述全部以此为准更正。满血 4090 比 4090D 强约 10\%。
\item 驱动 595.80,CUDA 13.2;功耗墙 450W,负载下 SM 时钟钉 2520MHz,无降频无虚拟化限速。
\item conda env \texttt{py312}(主栈):python 3.12.11 + torch 2.13.0+cu132 + transformers 5.5.4
+ av 15.1.0(pyav 后端;torchcodec 全系不兼容 torch 2.13)+ lerobot editable。
\item conda env \texttt{official}(官方对照栈):torch 2.8.0+cu128 + transformers 4.57.3 + flash-attn 2.8.3。
\item 网络:HF 直连 11.9MB/s $>$ hf-mirror;aliyun pypi;github 直连通;入口 18--30MB/s 共享。
\item 访问方式:本地 \texttt{/usr/bin/python3 bench/d4090.py exec|put|get}(paramiko,
密码文件 \texttt{\textasciitilde/.cache/.d4090\_pass})。
\end{itemize}

\subsection{远程目录地图}
\begin{table}[htbp]
\centering\footnotesize
\begin{tabular}{@{}lp{9.6cm}@{}}
\toprule
路径 & 内容 \\
\midrule
\texttt{/root/lerobot-port} & 我们的栈(pr3967-combined + E1 + vendored 修复),editable 安装 \\
\texttt{/root/ckpts/shakehands-lerobot} & 12GB 转换后 ckpt,最终 bake:eager attn + compile(predict\_velocity/prefix) + cuda preprocess + gridthw + cudagraph + num\_steps=7 \\
\texttt{/root/ckpts/lingbot-vla-v2-6b} & 35GB 原始上游 ckpt;\texttt{Qwen3-VL-4B-Instruct} tokenizer \\
\texttt{/root/datasets/rebot\_shakehands} & 100 段 / 73801 帧 / 30fps,7 维双臂 \\
\texttt{/root/d4090\_cfg} & 全部配置+bench 工具(见 §2.2)+ 各实验日志子目录 \\
\texttt{/root/lingbot-vla-v2-official} & 官方仓 clone;\texttt{/root/d4090\_official} 官方环境与其实验产物 \\
\texttt{/root/d4090\_rtc} & RTC 实测 agent 产物(三臂已完成,数字见 §4) \\
\bottomrule
\end{tabular}
\end{table}

\section{代码状态}

\subsection{分支与提交}
本地 \texttt{/home/nvidia/lerobot-pr3967} 分支 \texttt{pr3967-combined}:
\begin{itemize}
\item \texttt{517275d0} fork 合并(已推送到 miracle-techlink/lingbot\_vla\_v2\_lerobot main,
fast-forward 181 commits;推送用专用 deploy key + \texttt{GIT\_SS H\_COMMAND -F /dev/null},
见记忆文件 miracle-techlink-deploy-keys)。
\item \texttt{07835805} \textbf{E1}:去噪循环 while 条件(GPU 张量 $\to$ 每步 host sync 排水)
改预计算时间表 for 循环。eager -6.4\%(955.1$\to$893.6ms),compile 不变;两路径同 seed
输出\textbf{逐位一致}。
\item \texttt{0e1c67d8} vendored 修复:\texttt{preprcess\_grid\_thw} 原返回 \texttt{pos\_embeds=None}
导致缓存每次失效(且 CUDA graph capture 非法);改为返回真实张量,数值不变。
\item \texttt{df0cf53e} \textbf{路线 A 落地}:库级开关 \texttt{use\_cudagraph\_denoise}
(默认 False),去噪环整环 CUDA graph 捕获,四路逐位验证全零。详见 §5.2。
\item 后三个 commit 尚未推 fork(推送由 Will 发起)。
\end{itemize}

\subsection{bench 工具(全部参数化,新实验零代码)}
\begin{itemize}
\item \texttt{/root/d4090\_cfg/official\_style\_bench.py}:严格官方口径 bench。
参数:\texttt{<ckpt> [iters] [warmup] [mode] [step|whole] [extras]},
extras 支持 \texttt{nocompile / gridthw / cudnnbench / tritonmoe}(逗号组合)。
\item \texttt{profile\_ours.py} + \texttt{trace\_gap.py}:profiler + GPU 时间轴空隙分析。
\item \texttt{e1\_val.py} + \texttt{e1\_compare.py}:同 seed 数值验证模板(逐位对比),任何性能改动过一遍 10 分钟。
\item \texttt{cudagraph\_probe.py}:CUDA graph capture 探针(迭代到 v5)。
\end{itemize}

\section{今日实验台账(全部实测,严格口径 = cuda.Event 包住 sample\_actions)}

\subsection{2$\times$2 收官对比}
\begin{table}[htbp]
\centering\small
\caption{官方栈 vs 我们的栈(同机,满血 4090,独占 GPU)}
\begin{tabular}{@{}lrr@{}}
\toprule
 & 严格口径(纯前向) & 全链路 \\
\midrule
官方公开栈(FA2 + default compile) & 315.8ms & server 333 / client 338ms \\
\textbf{我们的栈(sdpa + max-autotune per-step)} & \textbf{235.0ms} & \textbf{261.6ms} \\
官方 README claim & 130ms(不可复现) & --- \\
\bottomrule
\end{tabular}
\end{table}

\subsection{实验记录(按时间序)}
\begin{table}[htbp]
\centering\footnotesize
\caption{全部 A/B 臂与结论(2026-08-14 基线 + 08-15 收官)}
\begin{tabular}{@{}p{4.8cm}p{2.6cm}p{6.6cm}@{}}
\toprule
实验臂 & 结果 & 结论 \\
\midrule
compile scope: per-step + max-autotune-no-cudagraphs(锁定) & \textbf{235.0ms} & 基线 \\
\quad whole 整图 + max-autotune & 252.1ms & 整图更慢(+7.3\%),否决 \\
\quad whole/step + default mode & 325.4 / 331.0ms & mode 比 scope 重要,锁定 max-autotune \\
\midrule
precompute\_grid\_thw=True & 233.7ms & 噪声级微赚,零风险,建议 bake \\
cudnn.benchmark=True & 238.4ms & 无收益,否决 \\
\midrule
\textbf{E1}(杀 while host-sync) & eager -61.5ms & \textbf{采纳}(07835805),逐位一致 \\
E1 对 compile 路径 & +0.5ms & 无变化 $\Rightarrow$ 步间 host 开销已被隐藏 \\
\midrule
\textbf{eager vs sdpa}(joint attn,4 步严格口径) & 145.3 vs 153.3ms & \textbf{采纳 eager}(-8.0ms):自定义 2D mask 下 sdpa 走 math 路径,eager 逐点 op 可被 inductor 融合 \\
\quad eager vs sdpa(graph e2e) & 107.5 vs 124.5ms & eager -17.0ms;ViT 内部仍 sdpa \\
\midrule
\textbf{路线 A}:CUDA graph 去噪环捕获 & \textbf{4 步 107.5ms e2e} & \textbf{打通}(df0cf53e),四路逐位验证全零,见 §5.2 \\
\midrule
Triton grouped-GEMM MoE 臂 & 723.6ms & \textbf{否决};bs=1 带宽受限,dense 双 GEMM 是正解 \\
\midrule
官方栈观测缩放(1cam256 $\to$ 3cam512) & 307$\to$386ms & 曲线近水平,外推 0 图像 token floor $\approx$300ms \\
官方栈 profiler 解剖 & GPU 忙时 113.7ms & launch 18,865 次,空隙 $\approx$200ms \\
我们的栈 profiler 解剖 & GPU 忙时 $\approx$172.6ms & launch 27,176 次,空隙 $\approx$60--75ms \\
\bottomrule
\end{tabular}
\end{table}

\subsection{官方 130ms 终审(两路独立调查,证据互锁)}
\begin{enumerate}
\item \textbf{claim 出处}:README 首 commit(2026-07-07)即有,论文/issues/bench 脚本零佐证;
附带的命令本身有误(\texttt{-{}-use\_compile} 裸旗标对 str2bool 直接报错),该节未逐字验证过。
\item \textbf{观测规模不是答案}:缩放曲线在官方配置区间是平的;不喂图像也有 $\approx$300ms floor。
\item \textbf{开关没有没开的}:FA2(ViT 侧)与官方 Triton MoE(robby\_moe)实锤生效;
joint attention 官方自研三路分发,deploy 主动选 eager(flex\_cached 推理即崩、flex 慢 3倍);
\texttt{seed\_kernels} 是字节私有依赖(注册代码被注释);\texttt{group\_gemm/kernel/} 只是
fallback 依赖,推理一次都没执行。
\item \textbf{数字解剖}:官方栈 GPU 纯忙时 \textbf{113.7ms $\approx$ 130ms}——claim 对应的是
``GPU 忙时''口径,只能出自内部 CUDA Graph 化构建;或注释掉的 \texttt{num\_steps=4}
(deploy:309,4 步 $\approx$150ms),但那与 README 自述的 10 步矛盾。
\item 硬件方向:我们的是满血 4090,比 claim 的 4090D 还强 $\approx$10\%,不利于 claim。
\end{enumerate}

\subsection{数值验证纪律}
E1 验证方式(模板化):同 seed 各 3 个 chunk,eager/compile 两路径 pre/post 对比。
E1 结果:两路径 \textbf{max$|\Delta|$=0,100\% 逐位一致}。后续任何性能改动都必须过这道。

\section{RTC / async 调研与实测结论(调度层)}
\begin{itemize}
\item 本分支自带完整 RTC(\texttt{policies/rtc/},PI 论文 arXiv 2506.07339 的 lerobot 实现)
与 gRPC async\_inference(老方案)。
\item \textbf{RTC 三臂实测完成}(30fps 仿真回路):等效反应延迟
\textbf{sync $\approx$1.0--1.2s $\to$ async $\approx$700ms $\to$ RTC $\approx$580ms}。
RTC 是推荐路线:本地线程无网络,引导式 chunk 融合,不改任何数值。
\item \textbf{引导开销仅 +9.1\%}——远低于此前预估的 +28\%。原因:lerobot 的
\texttt{RTCProcessor} 梯度路径是恒等映射,冻结参数后无额外 backward。
\item 两个适配坑:① 我们 \texttt{predict\_action\_chunk} 内部直接反归一化,RTC 要归一化空间
chunk,需拆分支;② 带梯度路径与 compile 的相互作用要重计时。
\item \textbf{LatencyTracker 坑}:\texttt{LatencyTracker.max()} 单调递增,长跑后读数失真,
应改 p95 窗口。
\item gRPC async(weighted-average 融合)不推荐:同机场景无理由选用,且论文显示延迟抖动下会振荡。
\end{itemize}

\section{后续优化路线:torch 原生手段消 launch 开销(不引自定义算子)}

\subsection{问题定位}
我们的 235ms = GPU 忙时 $\approx$172.6ms + 非 GPU 空隙 $\approx$60--75ms。
空隙来源:10 次 compiled 图进出(guard 评估+Python 胶水)、prefix 侧 graph break
残留、步间 eager 小张量 op。官方栈证明:同模型 GPU 忙时可到 113.7ms——
我们 GPU 侧还有 $\approx$59ms kernel 质量差距(他们的 FA2-ViT 与调优 MoE),
host 侧还有 60--75ms 空隙。\textbf{两边都有肉,都不需要自定义算子。}

\textbf{约束澄清}:``不用 GEMM/Triton 算子''指不引入\textit{手写/外部} kernel 依赖;
\texttt{torch.compile} 内部生成的 triton 与 \texttt{torch.cuda.CUDAGraph} 都是
\textbf{torch 原生 API},合规。社区版全部做成默认关的 flag。

\subsection{路线 A:CUDA graph 去噪环捕获(\textbf{已打通},commit \texttt{df0cf53e})}
\begin{itemize}
\item \textbf{落地形态}:库级开关 \texttt{use\_cudagraph\_denoise}(新增,默认 False)。
整个去噪环捕成单次 CUDA graph replay——27,176 次 launch 变 1 次。prefix 保持
compiled 不动,其输出 KV 经 \texttt{foreach\_copy\_} staging 进静态 buffer。
\item \textbf{关键实现要点}:① \textbf{空 cache 捕获}——只捕去噪环,完全绕开 ViT
非 FA2 分支的 \texttt{lengths.tolist()}(D2H sync,原 §5.2 的当前 blocker,已解);
② \texttt{foreach\_copy\_} staging 写静态 buffer;③ \texttt{fail\_on\_recompile}
防静默重编译;④ 观测 shape 变化自动回退普通循环(warning)。CUDA only;首次调用
多两次 warm+capture。
\item \textbf{精度}:四路逐位验证全零(变观测 / 首调用 / sdpa / eager / 有无 compile,
全部 max$|\Delta|$=0),E1 模板过录。
\end{itemize}

\begin{table}[htbp]
\centering\small
\caption{路线 A 步数扫描(RT$\!$X 4090,最终 bake 配置;probe 口径 = e2e,库级口径 = cuda.Event 包住 sample\_actions)}
\begin{tabular}{@{}rrrr@{}}
\toprule
步数 & e2e probe (ms) & 库级 cuda.Event (ms) & graph 关对照 (ms) \\
\midrule
4 & 107.5 & 108.5 & 135.5 \\
6 & 122.5 & --- & --- \\
7 & 130.1 & 130.9 & 175.9 \\
8 & 139.4 & --- & --- \\
10 & 153.6 & 153.0 & 227.0 \\
\bottomrule
\end{tabular}
\end{table}

graph 收益随步数放大:4/7/10 步分别 -27/-45/-74ms。7 步 130.9ms 与官方 README 的
130ms claim 同数字,但我们的公开可复现(官方 claim 不可复现,§3.3)。

\textbf{最终 bake 配置}(全部写进转换后 ckpt 的 \texttt{config.json},加载即用):
\begin{itemize}[itemsep=0pt,topsep=1pt,parsep=0pt]
\item \texttt{num\_steps=7}、\texttt{precompute\_grid\_thw=true}、
\texttt{use\_cudagraph\_denoise=true}
\item \texttt{attention\_implementation=eager}(joint attn,ViT 内部仍 sdpa)、\texttt{dtype=bf16}
\item \texttt{compile\_predict\_velocity=true}(per-step, max-autotune-no-cudagraphs)
\item \texttt{compile\_prefix=true}、\texttt{preprocess\_device=cuda}
\end{itemize}

\subsection{路线 B:无 cudagraph 的纯 compile 深挖(并行探)}
\begin{itemize}
\item \textbf{denoise-loop 单图 compile}(精确未测配置):只把 10 步循环编一张图,
prefix 独立——今天的 whole 臂(252ms)被 prefix 的 graph break 污染,不是公平测试。
预期回收 9/10 图进出开销,20--40ms。
\item \textbf{inductor 调优}(缩小 GPU 忙时那 59ms 差距):
\texttt{coordinate\_descent\_tuning=True}、\texttt{max\_autotune\_gemm\_backends} 加
CUTLASS(torch 自带,非外部依赖)、\texttt{freezing=True}。对 M=51 skinny GEMM 有
公开证据的个位数到 15\% kernel 收益。
\item prefix graph break 收尾:\texttt{precompute\_grid\_thw} 修复后(0e1c67d8)稳态
可全缓存,再 bake 进 ckpt。
\end{itemize}

\subsection{路线 C:调度层(与 A/B 正交,可叠加)}
RTC 移植(1--2 天)+ 真机 A/B。不改变单 chunk 时延,但有效反应延迟约减半
(实测 §4:sync 1.0--1.2s $\to$ RTC $\approx$580ms),是真机体验的最大单项。
三臂实测已完成,定稿参数:delay=8、execution\_horizon=10--16、废弃 use\_length。

\subsection{明确不做的}
手写/外部 Triton kernel(in-tree MoE 实测 723.6ms 否决)、flash-attn 包(sm\_89 可用但
官方实测证伪其价值)、FA3/FA4(架构不支持)、fp8/nvfp4 量化(撞精度红线)、
SageAttention(量化 attention 无机器人 drift 证据)、vLLM/SGLang fused MoE 借用
(bs=1 证据持平或更慢)、跨调用 prefix KV 精确复用(数学上不成立,图像变文本 KV 也变)。

\section{效率基建(今天沉淀的,下次直接用)}
\begin{itemize}
\item skill \texttt{remote-gpu-bringup}:新机器建链 playbook(测速选源/conda-pack 环境分发/
aria2c 断点续传/4-agent 分工纪律/ssh 坑位),目标 1 小时建链。
\item 实验套路:参数化 bench + driver 串行多臂 + 后台 monitor 自动收尾 + GPU 分区并行
(一卡一实验)+ 数值验证模板。今天 4 卡最多 3 实验并行零干扰。
\item ssh 纪律:nohup 会挂 exec 但任务已起(看 log 别重发);pgrep -f 自匹配;
sftp 大文件断点用 dd 续传。
\end{itemize}

\section*{在途事项(交接时状态)}
\begin{itemize}[itemsep=0pt,topsep=1pt,parsep=0pt]
\item RTC 三臂实测\textbf{已完成}:sync $\approx$1.0--1.2s / async $\approx$700ms /
RTC $\approx$580ms,引导 +9.1\%,数字已回填 §4;LatencyTracker 改 p95 窗口的坑已记录。
\item CUDA graph 去噪环 capture(路线 A)\textbf{已完成}:df0cf53e,步数扫描与验证见 §5.2。
\item fsdp2 训练\textbf{已打通}。2$\times$24GB 需四件套,缺一不可:
\texttt{train\_expert\_only=true}、auto\_wrap 三类层、
\texttt{fsdp\_offload\_params=true}、policy() bug 修复;$\approx$23s/步(offload 代价)。
单卡 24GB 走 gradient\_checkpointing,峰值 22.0GB,$\approx$0.7s/步。
全参微调 2$\times$24GB 物理不可行(96GB 状态/2 卡)。
\item E1、vendored 修复、use\_cudagraph\_denoise 三个 commit 待推 fork(等 Will 发起)。
\item 文档更正:pipeline PDF 等处的 ``4090D/sm\_120'' 应改为 ``4090/sm\_89''。
\end{itemize}

\end{document}
